% GENERATED FILE. Edit canonical lessons, then run npm run generate:books.
% canonical-source-hash: fnv1a64:f691f7a4257b28fa
% canonical-lessons: JA-C34-suika, JA-C34-daikon, JA-C34-ninjin, JA-C34-toufu, JA-C34-shio

\chapter{Watermelon, White radish, Carrot, Tofu, Salt}
\label{ch:ja-watermelon-white-radish-carrot-tofu-salt}
\hlchaptermodality{34}

\begin{chapteropening}
\textbf{By the end of this chapter:} \emph{I can say a watermelon, a white radish, a carrot, tofu and salt.}

Complete the last lesson of chapter 34: I can say a watermelon, a white radish, a carrot, tofu and salt.
\end{chapteropening}

\section[\texorpdfstring{\ja{すいか} (suika)}{すいか (suika)}]{\ja{すいか} (suika) — a watermelon}

\label{lesson:JA-C34-suika}



\begin{quote}
Before the new one: say the Japanese for a peach, then the Japanese for a pear.
\end{quote}

\subsection*{You'll want to know: \ja{すいか}}
\textbf{\ja{すいか}} — \emph{suika} — \textquotedblleft{}a watermelon\textquotedblright{}.

\begin{cousinweb}
\textbf{\ja{すいか}} (\emph{suika}) is Sino-Japanese, borrowed from Chinese, \textquotedblleft{}western melon\textquotedblright{}: \emph{sui}, \textquotedblleft{}west\textquotedblright{}, and \emph{ka}, \textquotedblleft{}melon\textquotedblright{} (Middle Chinese \emph{kwae}). \emph{Sui} is a later Chinese reading of the character for west, not the usual \emph{sei} (JMdict; Unihan, KANJIDIC2).
\end{cousinweb}

\begin{grammarlens}[title={What you've built}]
One more word for this chapter.
\end{grammarlens}

\subsection*{Your turn}
\begin{itemize}
\raggedright
  \item \emph{Say it:} \emph{suika}
  \item \emph{Say it:} \emph{suika}, once more
  \item \emph{Say it:} say \emph{suika}
  \item \emph{Recall:} say the Japanese for a peach, then the Japanese for a pear, then say \emph{suika} again
\end{itemize}

\begin{tcolorbox}[breakable,colback=teal!4,colframe=teal!35!black,title={Before you move on}]
What does \ja{すいか} mean? (\textquotedblleft{}A watermelon\textquotedblright{}.) Say it once more.
\end{tcolorbox}
\section[\texorpdfstring{\ja{だいこん} (daikon)}{だいこん (daikon)}]{\ja{だいこん} (daikon) — a white radish}

\label{lesson:JA-C34-daikon}



\begin{quote}
Before the new one: say the Japanese for a pear, then the Japanese for a watermelon.
\end{quote}

\subsection*{You'll want to know: \ja{だいこん}}
\textbf{\ja{だいこん}} — \emph{daikon} — \textquotedblleft{}a white radish\textquotedblright{}.

\begin{cousinweb}
\textbf{\ja{だいこん}} (\emph{daikon}) is Sino-Japanese (kango): two Chinese characters read with their on'yomi, readings borrowed from Chinese. They are \emph{dai}, \textquotedblleft{}big\textquotedblright{} (Middle Chinese \emph{dajH}) and \emph{kon}, \textquotedblleft{}root\textquotedblright{} (\emph{kon}) (Unihan, KANJIDIC2; Middle Chinese after Baxter–Sagart).
\end{cousinweb}

\begin{grammarlens}[title={What you've built}]
One more word for this chapter.
\end{grammarlens}

\subsection*{Your turn}
\begin{itemize}
\raggedright
  \item \emph{Say it:} \emph{daikon}
  \item \emph{Say it:} \emph{daikon}, once more
  \item \emph{Say it:} say \emph{daikon}
  \item \emph{Recall:} say the Japanese for a pear, then the Japanese for a watermelon, then say \emph{daikon} again
\end{itemize}

\begin{tcolorbox}[breakable,colback=teal!4,colframe=teal!35!black,title={Before you move on}]
What does \ja{だいこん} mean? (\textquotedblleft{}A white radish\textquotedblright{}.) Say it once more.
\end{tcolorbox}
\section[\texorpdfstring{\ja{にんじん} (ninjin)}{にんじん (ninjin)}]{\ja{にんじん} (ninjin) — a carrot}

\label{lesson:JA-C34-ninjin}



\begin{quote}
Before the new one: say the Japanese for a watermelon, then the Japanese for a white radish.
\end{quote}

\subsection*{You'll want to know: \ja{にんじん}}
\textbf{\ja{にんじん}} — \emph{ninjin} — \textquotedblleft{}a carrot\textquotedblright{}.

\begin{cousinweb}
\textbf{\ja{にんじん}} (\emph{ninjin}) is Sino-Japanese (kango): two Chinese characters read with their on'yomi, readings borrowed from Chinese. They are \emph{nin}, \textquotedblleft{}person\textquotedblright{} (Middle Chinese \emph{nyin}) and \emph{jin}, \textquotedblleft{}ginseng\textquotedblright{} (\emph{srim}) (Unihan, KANJIDIC2; Middle Chinese after Baxter–Sagart). JMdict lists ginseng as a second sense: the name first belonged to that root and later passed to the carrot.
\end{cousinweb}

\begin{grammarlens}[title={What you've built}]
One more word for this chapter.
\end{grammarlens}

\subsection*{Your turn}
\begin{itemize}
\raggedright
  \item \emph{Say it:} \emph{ninjin}
  \item \emph{Say it:} \emph{ninjin}, once more
  \item \emph{Say it:} say \emph{ninjin}
  \item \emph{Recall:} say the Japanese for a watermelon, then the Japanese for a white radish, then say \emph{ninjin} again
\end{itemize}

\begin{tcolorbox}[breakable,colback=teal!4,colframe=teal!35!black,title={Before you move on}]
What does \ja{にんじん} mean? (\textquotedblleft{}A carrot\textquotedblright{}.) Say it once more.
\end{tcolorbox}
\section[\texorpdfstring{\ja{とうふ} (toufu)}{とうふ (toufu)}]{\ja{とうふ} (toufu) — tofu}

\label{lesson:JA-C34-toufu}



\begin{quote}
Before the new one: say the Japanese for a white radish, then the Japanese for a carrot.
\end{quote}

\subsection*{You'll want to know: \ja{とうふ}}
\textbf{\ja{とうふ}} — \emph{toufu} — \textquotedblleft{}tofu\textquotedblright{}.

\begin{cousinweb}
\textbf{\ja{とうふ}} (\emph{toufu}) is Sino-Japanese (kango): two Chinese characters read with their on'yomi, readings borrowed from Chinese. They are \emph{tō}, \textquotedblleft{}bean\textquotedblright{} (Middle Chinese \emph{duwH}) and \emph{fu}, \textquotedblleft{}rot, curdle\textquotedblright{} (\emph{bjuX}) (Unihan, KANJIDIC2; Middle Chinese after Baxter–Sagart).
\end{cousinweb}

\begin{grammarlens}[title={What you've built}]
One more word for this chapter.
\end{grammarlens}

\subsection*{Your turn}
\begin{itemize}
\raggedright
  \item \emph{Say it:} \emph{toufu}
  \item \emph{Say it:} \emph{toufu}, once more
  \item \emph{Say it:} say \emph{toufu}
  \item \emph{Recall:} say the Japanese for a white radish, then the Japanese for a carrot, then say \emph{toufu} again
\end{itemize}

\begin{tcolorbox}[breakable,colback=teal!4,colframe=teal!35!black,title={Before you move on}]
What does \ja{とうふ} mean? (\textquotedblleft{}Tofu\textquotedblright{}.) Say it once more.
\end{tcolorbox}
\section[\texorpdfstring{\ja{しお} (shio)}{しお (shio)}]{\ja{しお} (shio) — salt}

\label{lesson:JA-C34-shio}



\begin{quote}
Before the new one: say the Japanese for a carrot, then the Japanese for tofu.
\end{quote}

\subsection*{You'll want to know: \ja{しお}}
\textbf{\ja{しお}} — \emph{shio} — \textquotedblleft{}salt\textquotedblright{}.

\begin{cousinweb}
\textbf{\ja{しお}} (\emph{shio}) is a native Japanese word, inherited rather than borrowed from Chinese, and its older origin is not settled. Its kanji was attached to it later as a kun'yomi, a native reading given to a Chinese character (JMdict, KANJIDIC2).
\end{cousinweb}

\begin{grammarlens}[title={What you've built}]
One more word for this chapter.
\end{grammarlens}

\subsection*{Your turn}
\begin{itemize}
\raggedright
  \item \emph{Say it:} \emph{shio}
  \item \emph{Say it:} \emph{shio}, once more
  \item \emph{Say it:} say \emph{shio}
  \item \emph{Recall:} say the Japanese for a carrot, then the Japanese for tofu, then say \emph{shio} again
\end{itemize}

\begin{tcolorbox}[breakable,colback=teal!4,colframe=teal!35!black,title={Before you move on}]
What does \ja{しお} mean? (\textquotedblleft{}Salt\textquotedblright{}.) Say it once more.
\end{tcolorbox}
